% Short running form: the full title is wider than the box the running head
% leaves beside the Sunday's name.
\runninghead{The King's Feast for All Peoples}
\properlane{wedding-feast}{The King's Feast for All Peoples: The Wedding of the Son and the Banquet on the Mountain}
\label{sec:the-kings-feast}

Read from its centre, the Sunday is about what God gives. A king makes a
wedding for his son and will not let the hall stand empty; the Lord of hosts
makes a feast on his mountain for all peoples and swallows up death there. The
two titles the \work{Ordo} sets over these readings put them side by side. The
witnesses speak each of his own passage. Gregory the Great reads the parable's
wedding as the Church joined to the Son, and places it in the present Church;
Hilary of Poitiers reads it as the sacrament of the glory to be received at
the resurrection; Jerome reads Isaiah's banquet as the feast the Lord makes
for the nations after the Passion; and Aquinas alone sets Isaiah's feast
beside the king's dinner. This interpretation reads the two feasts as one
gift: the wedding the Father makes for his Son, to which the call goes out to
everyone the servants can find, spread now in the Church and completed in the
resurrection. That reading of the whole is the editor's synthesis, which no
witness cited holds as stated, and so is the joining of both readings to this
Mass's psalm, prayers and Communion.

\subsection{A king who made a marriage for his son}

The parable's first sentence already raises the question the Fathers answer
first. Who is the king, who is the son, and what is the wedding? Gregory the
Great begins from the king: he is the one to whom the Psalmist says,
\latin{Deus iudicium tuum regi da, et iustitiam tuam filio regis}, give to the
king thy judgment, O God, and thy justice to the king's son (Ps 71:2). Then he
says when the wedding was made, and corrects himself as he says it. God the
Father made the wedding for God his Son, he says first, \latin{quando hunc in
utero Virginis humanae naturae coniunxit}, when in the Virgin's womb he joined
him to human nature. But a joining of that kind is ordinarily made of two
persons, and he will not have his hearers think the person of the Redeemer
made out of two: \latin{absit hoc ab intellectibus nostris}. So he restates
it: \latin{Apertius ergo atque securius dici potest quia in hoc Pater regi
Filio nuptias fecit, quo ei per incarnationis mysterium sanctam Ecclesiam
sociavit}, it can be said more openly, then, and more safely, that the Father
made the wedding for the King his Son in this, that through the mystery of the
Incarnation he joined the holy Church to him. The Virgin's womb was the
bridegroom's chamber, out of which he came forth ``as a bridegroom'' (Ps 18:6).
The servants sent twice are the prophets, who said that the Incarnation would
come, and the apostles, who announced that it had come (\work{Homiliae in
Evangelia} 38.3).

Jerome states the same identification in two sentences and adds who the Church
is: \latin{Rex iste qui fecit nuptias filio suo, Deus omnipotens est. Facit
autem nuptias Domino nostro Iesu Christo et Ecclesiae, quae tam ex Iudaeis,
quam ex gentibus, congregata est}: this king who made a wedding for his son is
God almighty, and he makes the wedding for our Lord Jesus Christ and for the
Church, which is gathered as much from the Jews as from the gentiles
(\work{Commentarii in Matthaeum} III, PL 26, col.~159). Irenaeus reads the
parable against teachers who divided the God who calls from the God who
judges. The Lord shows by it, he says, ``that there is one King and Lord, the
Father of all'', who ``had from the beginning prepared the marriage for His
Son'', sent his servants to call the invited, and, when they would not come,
``called together from all the highways, that is, from all nations, [guests] to
the marriage feast of His Son''. The same King ``grants them the incorruptible
banquet'' (\work{Against Heresies} IV.36.5--6). For Irenaeus the point is the
unity of the design: one Father prepares the wedding from the beginning,
issues the first call and the second, and fills the hall.

John Chrysostom asks why the parable speaks of a marriage at all, and answers
from what a wedding is: ``That thou mightest learn God's tender care, His
yearning towards us, the cheerfulness of the state of things, that there is
nothing sorrowful there, nor sad, but all things are full of spiritual joy.''
He sets beside it Paul's ``I have espoused you to one husband'' and ``This is a
great mystery, but I speak concerning Christ and the Church'', and he adds that
the parable, coming after the vineyard and the slain son, speaks of a wedding
after a death: ``Hereby He proclaimed the resurrection also'', for ``even after
the death, then is the marriage, then the bridegroom'' (\work{Homilies on
Matthew} 69.1). Hilary of Poitiers draws that last thread tightest. The
wedding, he says, is \latin{vitae coelestis et in resurrectione suscipiendae
aeternae gloriae sacramentum}, the sacrament of the heavenly life and of the
eternal glory to be received in the resurrection; and it is rightly made by
the Father, \latin{quia aeternitatis huius societas, et novi corporis
desponsa coniunctio, iam perfecta habebatur in Christo}, because the fellowship
of this eternity and the espoused union of the new body were already held
perfect in Christ (\work{Commentarius in Matthaeum} 22.3, PL 9, col.~1042).

These witnesses agree that the king is the Father and the son is Christ, that
the wedding is the Father's gift to the Son, and that the invitation goes out
beyond those first invited. They do not say the same thing
about the bride. Gregory names the Church, joined to the Son through the
Incarnation; Jerome the Church of Jews and gentiles; Irenaeus speaks of the
marriage of the Son and of the guests from all nations; Hilary speaks of the
new body already united in Christ and of the glory it will receive.

\subsection{``My dinner is prepared''}

The second invitation describes the feast: the dinner is ready, the beeves and
fatlings are killed, all things are ready (Mt 22:4). Chrysostom hears in it
what the guests are being called to: ``Unto labors, and toils, and sweat? Nay
but unto pleasure. \ldots\ See how complete His banquet, how great His
munificence.'' Jerome reads the menu as a figure: \latin{per metaphoram opes
regiae describuntur, ut ex carnalibus intelligantur spiritualia; vel certe
dogmatum magnitudo}, the king's riches are described by metaphor, so that
spiritual things may be understood from bodily ones, or else the greatness of
the doctrines (col.~159). Gregory reads the slain animals as people. The bulls
are the fathers of the old covenant, who struck their enemies as with a horn;
the fatlings are the fathers of the new, who \latin{gratiam pinguedinis
internae percipiunt}, receive the grace of an inner fatness, and are lifted by
it to contemplation, as the Psalmist desired, ``Let my soul be filled as with
marrow and fatness'' (Ps 62:6). That they are named as slain only at the
second invitation he takes as God's way of adding examples to words when words
are not heard: \latin{Patrum praecedentium mortes aspicite, et remedia vitae
vestrae cogitate}, look on the deaths of the fathers who went before, and think
of the remedies for your own life (38.4). Hilary too reads the fatlings as
spiritual men, fed \latin{tamquam coelesti pane}, as if with heavenly bread
(22.4, col.~1043).

The table is not only described; it is set. Gregory says to his own hearers,
in the homily's application, \latin{Nos sumus, fratres charissimi, qui in
nuptiis Verbi discumbimus, qui iam fidem in Ecclesia habemus, qui Scripturae
sacrae epulis pascimur, qui coniunctam Deo Ecclesiam esse gaudemus}: we are
the ones, dearest brethren, who recline at the wedding of the Word, who
already have faith in the Church, who are fed with the dishes of holy
Scripture, who rejoice that the Church is joined to God (38.11). The feast is
present, and he names its food as the Scriptures read in the assembly.

\subsection{On this mountain, for all peoples}

Isaiah's feast is the other half of the official correlation, and two
witnesses read it of Christ. Jerome's sentence carries a contrast that is his
own and not the prophet's: \latin{Post passionem ergo Domini \ldots\ faciet
Dominus nequaquam populo Iudaeorum, sed omnibus gentibus in monte Sion pingue
convivium}, after the Lord's Passion the Lord will make, not for the people of
the Jews but for all nations, a rich feast on Mount Sion, so as to cast down
and swallow up ``the face of death'' and of the bond by which all peoples were
bound. He reads the swallowing of death with Paul, \latin{iuxta Apostolum,
absorbebitur mors in perpetuum}, and the wiping of tears as the moment
\latin{quando morte superata, Christi advenerit regnum}, when death is overcome
and Christ's kingdom has come (\work{Commentarii in Isaiam} VIII, PL 24,
col.~300). The verse he expounds says ``unto all people'', and the same
oracle promises that the reproach of the Lord's own people will be taken away
from the whole earth (Is 25:8); the exclusion in Jerome's sentence is his
reading, which the last section of this interpretation weighs.

Thomas Aquinas makes the join between the two readings himself. The feast is
made on this mountain \latin{quia ibi Christus passus est, unde nobis omnia bona
provenerunt; vel quia ibi iudicabit, ut dicitur Matth. 22: Ecce prandium meum
paravi, tauri mei et altilia sunt occisa}: because Christ suffered there, from
which all good things have come to us, or because he will judge there, as
Matthew 22 says, ``Behold, I have prepared my dinner; my bulls and fatlings are
killed''. In a note on the feast of fat things he then distinguishes three
feasts. The first is \latin{familiare militantis Ecclesiae}, the household
feast of the Church on pilgrimage, for which he cites Paul's ``as often as you
shall eat this bread and drink the chalice, you shall shew the death of the
Lord, until he come'' (1~Cor 11:26) and this Sunday's psalm, \latin{Impinguasti
in oleo caput meum}. The second is \latin{convivium privatum animae}, the
private feast of the soul. The third is \latin{convivium solemne caelestis
curiae}, the solemn feast of the heavenly court, whose sweetness is given
\latin{ad satietatem}, to the full, with the psalm verse \latin{Satiabor cum
apparuerit gloria tua}, ``I shall be satisfied when thy glory shall appear''
(Ps 16:15) (\work{Expositio super Isaiam} c.~25, Parma, pp.~501--502). So
Aquinas holds together the Gospel's dinner, Isaiah's mountain, the psalm's
anointed head and the altar, and he does it as an exegete of Isaiah; it is the
one documented joining of this Sunday's readings among the witnesses cited.

The psalm's table receives an analogous threefold reading in Aquinas's
commentary on the Psalter, where the table is Scripture's teaching \latin{vel
mensam sacramentalem, scilicet altaris}, and the three tables are those of the
old law, of the New Testament, and of the homeland (\work{Postilla super
Psalmos} 22); the triad is not the Isaiah commentary's, though both end in
the heavenly feast. Augustine reads the table as the solid food given to the grown, and the
mercy that follows as leading ``that I may dwell in the house of the Lord for
ever''; Bellarmine finds the Eucharist in the pasture; the fourth
\work{Mystagogical Catechesis} calls the table ``that mystical and spiritual
Table''. That the feast is the altar's is the reading of Aquinas and the
\work{Catechesis}; Isaiah and the psalmist do not say so, and the literal
sense of both is a banquet God gives.

\subsection{Now, and at the resurrection}

Where does the feast stand? Here the witnesses differ in emphasis, and the
difference is worth keeping. Gregory places the parable's wedding in the
present Church. The mixture of the guests shows it: \latin{per has regis
nuptias praesens Ecclesia designatur, in qua cum bonis et mali conveniunt},
the present Church is signified by this king's wedding, in which the bad come
together with the good (38.7); and Luke's supper, from which no one is cast
out, is for him the eternal and last banquet (38.1). Augustine speaks of two
feasts: ``there are two feasts of the Lord; one to which the good and evil
come, the other to which the evil come not''; and he joins them in a single
exhortation, ``Be that feast which now is, received worthily, that we may
attain to the other'' (\work{Sermo} 90.1, 4). Hilary makes the wedding the
sacrament of the glory to be received at the resurrection, already perfect in
Christ. Aquinas gives four expositions of the wedding, among them the Word
incarnate as bridegroom and the Church as bride (Eph 5:32) and the wedding
\latin{in communi resurrectione \ldots\ quando mortale nostrum absorbebitur a
vita}, in the common resurrection, when what is mortal in us will be swallowed
up by life; but for the parable he follows Gregory: \latin{secundum Gregorium,
oportet exponere de praesentibus}, according to Gregory it must be expounded of
present things (\work{Super Matthaeum} c.~22, Venice, p.~280). Chrysostom, as
above, sees the resurrection proclaimed in the wedding after the death.

These are emphases and not contradictions. Gregory does not read the
parable's wedding as the last banquet, and Hilary does not read it as the
Church on earth. What joins the present feast to the last differs with the
form of the Gospel. In both forms the link is the resurrection. Hilary's
wedding is the glory to be received in the resurrection, already perfect in
Christ, and Chrysostom hears the resurrection in a wedding made after a death.
Isaiah's verse supplies the end of the line. Jerome reads its swallowing of
death with Paul, and Aquinas reads the wedding of the resurrection in Paul's
words from another letter (2~Cor 5:4), so that ``Death is swallowed up in
victory'' (1~Cor 15:54), Isaiah's death cast down, and the mortal swallowed up
by life meet in their pages. Where the longer form is read, a second link
stands between the two feasts: every witness who comments on the king's
coming in to see the guests (22:11) places it at the judgment, and the
present feast becomes the anteroom of the last. Where the shorter form is
read, the king does not come in, and this interpretation keeps the two feasts,
the one spread now and the one completed in the resurrection, without the
judgment between them.

\subsection{The Missal's texts in this reading}

The Missal's antiphons and prayers serve all three years of the Lectionary,
and no witness cited connects them to these readings. In this interpretation
they take their places along the line of the feast. The Entrance Antiphon asks who could stand if the Lord marked
iniquities, and answers that with him there is propitiation. Augustine's
reading of the word, given with the antiphon above, makes the propitiation
the sacrifice offered for us, and so describes the door by which anyone enters
the feast: not his own standing but the sacrifice already made.
The Collect asks for grace that goes before and follows. In the sentence of
Augustine that Aquinas quotes, grace goes before \latin{ut vocemur}, that we
may be called, and follows \latin{ut glorificemur}, that we may be glorified;
heard in this reading, the call is the king's invitation and the glory is the
feast completed.

The second reading enters this interpretation through its last verses. In
the Vulgate, which the Douay--Rheims follows, v.~19 is a wish that God fill all
the Philippians' desire, \latin{impleat omne desiderium}; the \work{Nova
Vulgata} has the same \latin{desiderium} with a promise, \latin{implebit}.
Aquinas, after reading the verse as God's return to the givers, reads it of
glory: \latin{Et hoc verum in gloria \ldots\ quia ibi
implebitur totum desiderium}, and this is true in glory, because there all
desire will be filled, with the same psalm verse he gave to Isaiah's third
feast, \latin{Satiabor cum apparuerit gloria tua} (\work{Super Philippenses}
c.~4, lect.~2, p.~403). Joined to the feast, the verse rests on the Latin
\latin{desiderium}, desire, which both Latin texts read and the Douay--Rheims
renders ``want''. The Alleluia verse prays to know the hope of our calling,
and Chrysostom's comment on Paul's prayer, ``It is as yet, he means, hidden,
but not so to the faithful'', names the condition of those who have heard the
invitation and have not yet seen the feast. The Prayer over the Offerings asks
that through these acts of devout worship the faithful may pass to heavenly
glory: in this reading, the passage from the feast now spread to the feast
completed.

At Communion the Missal offers two antiphons, and this interpretation hears
each differently. Where the first is chosen, the rich who hungered are,
in Augustine's own joining of the verse to the Gospel of John, those who lack
the Living Bread; his words are quoted under the third interpretation. Where the second is chosen, ``we shall see him as he is''
answers the parable's king in both forms of the Gospel. Aquinas's sentence on
the king of its second verse, \latin{cum videbitur sicuti est, tunc erit Rex},
has the words of this verse, though he cites Paul's ``face to face''
(1~Cor 13:12) for it; Augustine's reading of the vision is given with the
antiphon above. Where the longer form is also
read, the antiphon answers as well the king's coming in to see the guests
(22:11): the guests see the King whose coming in is, for these witnesses, the
judgment. The Prayer after Communion asks that those fed with the
Body and Blood of Christ be made partakers of the divine nature. Heard in this
interpretation, it asks that the wedding Gregory describes in his more careful
formulation, the Church joined to the Son through the mystery of the
Incarnation, be shared by those fed at the table; the attributed
\work{Catechesis} joins the Body and Blood to Peter's ``partakers of the divine
nature'' (4.3), and Aquinas calls grace a participation of the divine nature.
It is Gregory's restated formulation that bears this reading; his first, of
the Son joined to human nature in the Virgin's womb, he corrected himself, as
shown above.

\subsection{Where this interpretation strains}

The parable is grim where Isaiah is glad. Isaiah's feast ends with tears wiped
from every face; the parable's ends, in the longer form, with weeping and
gnashing of teeth, and in both forms it passes through the slaughter of the
servants, the king's anger, the armies and the burnt city (Mt 22:6--7). An
interpretation that rejoices in the feast must carry that verse.

The Fathers read it in two ways. John Chrysostom says the parable foretells
``the casting out of the Jews, and the calling of the Gentiles'', and reads the
burning of the city as ``the things that took place under Vespasian and
Titus'', forty years after the Passion, ``that He might show His long
suffering'' (\work{Homilies on Matthew} 69.1). Jerome gives two readings of the
armies, avenging angels or the Romans under Vespasian and Titus (col.~160), and
Aquinas likewise gives both, adding that God's anger \latin{non commotionem
significat, sed vindictam}, signifies not agitation but just punishment
(p.~281). Gregory and Hilary read the armies as angels and the burning as the
eternal punishment of the persecutors: for Gregory \latin{illorum non solum
animae, sed et caro \ldots\ aeterna gehennae flamma cruciatur}, not only their
souls but their flesh is tormented by the eternal fire (38.5). Irenaeus,
Hilary, Chrysostom, Jerome and Aquinas identify the first invited with Israel,
Irenaeus as ``the men of the former dispensation''; Aquinas heads the parable
\latin{de reprobatione Iudaeorum, et vocatione gentium}; and Jerome's sentence on
Isaiah sets the feast \latin{nequaquam populo Iudaeorum}. Gregory identifies
the refusers with the persecutors of the preachers and applies the parable to
his own hearers.

These are the Fathers' readings of a parable that Matthew says was spoken to
the chief priests and the elders of the people (21:23) and that the chief
priests and Pharisees knew was spoken of them (21:45). The Second Vatican
Council, in \work{Nostra aetate} 4, governs how such readings are presented and
applied, though it passes no judgment on any Father's exegesis. It recalls
with Paul that the gifts and the calling of God are without repentance
(Rom 11:29; Latin, \latin{cuius dona et vocatio sine paenitentia sunt}); it
teaches that what happened in the Passion cannot be charged ``against all the
Jews, without distinction, then alive, nor against the Jews of today'', and
that the Jews ``should not be presented as rejected or accursed by God, as if
this followed from the Holy Scriptures''. A parable about a call is the place
where that bound bites hardest. Read under it, the identification of the
first invited belongs to the Fathers who make it; the parable's own hearers
are those leaders in the temple; and the warning falls, as Gregory and
Chrysostom themselves turn it, on the baptized. Gregory's ``We are the ones
who recline at the wedding of the Word'' and Chrysostom's ``Hear ye, as many as
having partaken of the mysteries'' (69.2) are addressed to Christians. Isaiah's
``all people'' cannot be narrowed to ``not this people'' on the strength of
Jerome's phrase, which is his alone.

Two further limits stand. The Eucharistic reading of Isaiah's feast and of the
psalm's table is the reception of Aquinas and of the attributed
\work{Catechesis}; the literal sense of both is a banquet God gives on his
mountain and in his house. And Gregory's reading of the wedding as the present
Church, which Aquinas follows, limits how far the parable's wedding may be
called the last banquet: the present feast and the final one are joined here
by the resurrection, and in the longer form also by the judgment, not
identified.

\subsection{The four senses of this interpretation}

\begin{fourSenses}
\item[Literal.] Isaiah promises a feast for all peoples on the Lord's mountain,
where death is cast down and tears are wiped away. Jesus tells the chief
priests and elders of a king whose wedding the invited refuse, and whose hall
he fills with whomever his servants find on the roads.

\item[Allegorical.] The king is the Father and the wedding is the Son's union
with the Church through the mystery of the Incarnation (Gregory); the Church
is gathered from Jews and gentiles alike (Jerome). Aquinas reads Isaiah's
feast on the mountain of the Passion as the feast announced by ``my dinner is
prepared'', and the psalm's table as the altar.

\item[Moral.] Come to the feast now spread, and receive it worthily, that you
may come to the other (Augustine). Where the first Communion antiphon is
chosen: hunger for the Living Bread, not for what leaves the rich empty
(Augustine).

\item[Anagogical.] The wedding is completed in the resurrection (Hilary), when
death is swallowed up and Christ's kingdom comes (Jerome, on Isaiah 25:8), and
the King is seen as he is (Aquinas). Where the second Communion antiphon is
chosen: we shall be like him, because we shall see him as he is.
\end{fourSenses}
